\documentclass[10pt,a4paper]{amsart}
\usepackage[margin=19mm]{geometry}
\usepackage{amsmath,amssymb,mathtools}
\usepackage[scr=rsfs,cal=euler]{mathalfa}
\usepackage{xfrac}
\usepackage[unicode,hidelinks,bookmarks=false]{hyperref}
\newcommand{\ie}{i.e.\ }
\newcommand{\cf}{cf.\ }
\newcommand{\resp}{respectively\ }
\newcommand{\Subsection}{\subsection}
\providecommand{\nobreakdash}{\nobreak\hbox{-}}
\setlength{\emergencystretch}{2em}
\multlinegap0pt
\usepackage{amssymb}
\usepackage[shortlabels]{enumitem}
\usepackage{mathtools}
\usepackage{xcolor}
\newtheorem{assumption}{Assumption}
\newtheorem{lemma}{Lemma}
\newcommand\R{{\mathbb{R}}}
\newcommand\dom{{\text{dom\,}}}
\newcommand\st{{\text{s.t.}}}
\newcommand\vA{{\vector{A}}}
\newcommand\vb{{\vector{b}}}
\renewcommand\vector[1]{\boldsymbol{#1}}
\newcommand\vf{{\vector{f}}}
\newcommand\vh{{\vector{h}}}
\newcommand\vx{{\vector{x}}}
\newcommand\vy{{\vector{y}}}
\newcommand\vzero{{\vector{0}}}
\title{Reachability-Augmented Dual Dynamic Programming for Optimal Path Parameterization}
\author{Yunan Wang, Jizhou Yan, Chuxiong Hu, Zeyang Li}
\date{Source: arXiv:2605.19089v1 (CC BY 4.0)}
\begin{document}
\maketitle

\noindent\emph{Source and licence.} Reachability-Augmented Dual Dynamic Programming for Optimal Path Parameterization. Version arXiv:2605.19089v1 (CC BY 4.0), distributed by arXiv under the Creative Commons Attribution 4.0 International licence, http://creativecommons.org/licenses/by/4.0/, as recorded in the arXiv licence metadata for this exact version and shown on the abstract page https://arxiv.org/abs/2605.19089v1. Full licence notice: \texttt{licenses/arxiv-2605.19089-CC-BY-4.0.txt}.\par\smallskip

\noindent\textit{Editorial scope.} Counted unit: the article's lemma lemma:convexity\_reachset\_valfunc (source0.tex lines 406--408). The article's complete original proof of it is delivered with it (source0.tex lines 1003--1017, one \texttt{proof} environment of 15 lines, which the article places away from the statement). No mathematics is written, repaired or re-derived: the statement and the proof are the article's own lines, in the article's own notation, and no retained line is substituted or re-flowed.  Retained uncounted as the source-local context the proof needs: lines 262--268; lines 286--288; lines 995--997.  Nothing was omitted from any retained span, so the package carries no omission record.  No retained line contains a citation, so the wrapper carries no bibliography and no bibliography entry.  No retained line contains a figure environment, an image-inclusion command or drawing code, so no image is delivered and the package declares no asset dependency.  Editorial formatting only, in addition: an AMS \texttt{amsart} preamble that restates the article's own declarations and macros used by the retained lines, namely the environments \texttt{assumption}, \texttt{lemma} on separate counters, together with the standard packages those lines need.  The retained environments print in this document as Assumption 1, Lemma 1 in order of appearance, whereas the article prints the same items with the numbering of its own full document.  The complete native source of the article as distributed in the arXiv e-print for this exact version is bundled unchanged as \texttt{sources/200892/source0.tex}.
    \begin{align}
        \min_{\vx_k,u_k} \quad &J = \sum_{k=0}^{N-1} L_k(\vx_k, u_k) + \Phi(\vx_N) \label{eq:op_discretize_objective}\\
        \st \quad & \vx_{k+1} = \vA_k\vx_k + \vb_k u_k, \,\,k=0,1,\ldots,N-1,\label{eq:op_discretize_dynamics} \\
        &\vf_k(\vx_k) \leq \vzero,  \,\,k=1,2,\ldots,N-1, \label{eq:op_discretize_state_constraints}\\
        &\vh_k(\vx_k, u_k) \leq \vzero, \,\,k=0,1,\ldots,N-1,\label{eq:op_discretize_mix_constraints}\\
        &\vx_0\in\mathcal{X}_0, \,\,\vx_N\in\mathcal{X}_\text{f}.
    \end{align}

\begin{assumption}\label{assum:feasibility_smoothness}
    Functions $L_k$, $\Phi$, $\vf_k$, and $\vh_k$ are locally Lipschitz continuous in their arguments for each $k$. The sets $\mathcal{X}_0$, $\mathcal{X}_\text{f}$, and $\{(\vx_k,u_k)\mid \vf_k(\vx_k)\leq\vzero,\,\vh_k(\vx_k,u_k)\leq\vzero\}$ are compact for each $k$.
\end{assumption}

\begin{lemma}\label{lemma:min_is_convex}
    Consider a convex function $V:\Omega\subset\R^m\times\R^n\to\R$ where $\Omega$ is a non-empty convex set. Define the function $U(\vx)=\min_{\vy\in\mathcal{Y}(\vx)} V(\vx,\vy)$, where the domain of $U$ is $\dom U=\{\vx\in\R^m\mid \mathcal{Y}(\vx)\neq\varnothing\}$, and $\mathcal{Y}(\vx)=\{\vy\in\R^n\mid (\vx,\vy)\in\Omega\}$. Then, $U$ is convex in the convex set $\dom U$.
\end{lemma}

\begin{lemma}\label{lemma:convexity_reachset_valfunc}
    In COPP, for each $k$, the backward reachable set $\mathcal{B}_k$ is convex and compact, whereas the value function $V_k^*$ is convex in $\mathcal{B}_k$. The admissible control set $\mathcal{C}_k(\vx_k)\subset\R$ is convex and compact for each $\vx_k\in\mathcal{B}_k$; hence, $\mathcal{C}_k(\vx_k)$ is either a closed interval or a singleton. Specifically, if \eqref{eq:op_discretize_state_constraints} and \eqref{eq:op_discretize_mix_constraints} are linear inequalities, then $\mathcal{B}_k$ is a polytope.
\end{lemma}

\begin{proof}[Proof of Lemma \ref{lemma:convexity_reachset_valfunc}]
    Consider the following convex set
    \begin{align}
        \widehat{\mathcal{B}}_k=\Big\{&(\vx_{k}, \ldots, \vx_N, u_k, \ldots, u_{N-1})\mid \notag\\
        &\forall k\leq j<N,\,\vx_{j+1} = \vA_j\vx_j+\vb_j u_j, \notag\\
        &\vf_j(\vx_j) \leq \vzero,\,\vh_j(\vx_j, u_j) \leq \vzero,\,\vx_N \in \mathcal{X}_\text{f} \Big\}.\label{eq:widehat_backward_set}
    \end{align}
    By Assumption \ref{assum:feasibility_smoothness}, $\widehat{\mathcal{B}}_k$ is compact. As a projection of $\widehat{\mathcal{B}}_k$ onto the $\vx_k$-subspace, $\mathcal{B}_k$ is also convex and compact. By Lemma \ref{lemma:min_is_convex}, $V_k^*$ is convex in $\mathcal{B}_k$ since $\widehat{V}_k:\widehat{\mathcal{B}}_k\to\R$ is convex:
    \begin{equation}
        \widehat{V}_k\left((\vx_j)_{j=k}^N, (u_j)_{j=k}^{N-1}\right) = \sum_{j=k}^{N-1} L_j(\vx_j, u_j) + \Phi(\vx_N).
    \end{equation}
    The convexity and compactness of $\mathcal{C}_k(\vx_k)$ is straightforward.

    If all the constraints are linear inequalities, then $\widehat{\mathcal{B}}_k$ is a polytope, and thus $\mathcal{B}_k$ is also a polytope.
\end{proof}

\end{document}
